\section{Learning framework}

Before diving into any algorithm, let us informally describe what we mean by
\emph{learning} from data. The guiding idea of machine learning is to replace
explicit, hand-written rules by a procedure that \emph{extracts patterns from
examples}. Rather than telling a computer exactly how to recognise a cat, or how
to price a flat, we provide it with a collection of \emph{data} and let it
\emph{learn} a rule that performs well and, crucially, keeps performing
well on new data it has never seen. This section sets up the vocabulary; the
precise definitions and the mathematics come in the following sections.

\subsection{What is ``data''?}

Everything starts from \emph{data}: a collection of examples
$x_1, \dots, x_n$, where $n$ is the number of samples. Throughout these notes we
represent each example as a vector (sometimes called ``sample'' or ``input'') $x_i$, but this vector can encode
essentially anything, provided we agree on a way to turn the raw object into
numbers. A few concrete cases:
\begin{itemize}
    \item \emph{Feature vectors.} The most direct case, where each coordinate is
    an interpretable measurement. To describe a flat, one might take
    \[
    x = \big( [\text{surface in m}^2],\ [\text{number of rooms}],\ [\text{floor}],\ \dots \big) .
    \]
    \item \emph{Images.} A colour image of $H \times W$ pixels is a grid of RGB
    intensities. Flattening it gives a vector $x \in \R^d$ with $d = 3 H W$ (one
    coordinate per colour channel and pixel). A modest $224 \times 224$ image
    already lives in dimension $d \approx 1.5 \times 10^5$.
    \item \emph{Sound.} An audio clip is an amplitude sampled through time. If we
    record $d$ amplitudes, the clip is a vector $x \in \R^d$, one coordinate per time step.
    \item \emph{Text, graphs, \dots} Similarly, a piece of text, a molecule, or a
    social network can all be encoded as vectors once we fix a suitable
    representation.
\end{itemize}
The common thread is that, after a suitable encoding, every input becomes a
point $x_i \in \R^d$. This is what lets a single mathematical framework cover
such different-looking objects.

In \emph{supervised} problems each input $x_i$ comes with a \emph{label} $y_i$,
and the nature of $y$ varies just as much. It may be a real number
$y \in \R$ (a price, a temperature), a binary label $y \in \{0, 1\}$ (spam or
not), one of several categories $y \in \{1, \dots, K\}$ (which digit is written),
or even a structured object such as another image or a sentence. The two basic
regimes, real-valued and categorical labels, are exactly the regression and
classification problems we distinguish below.

\subsection{Supervised vs.\ unsupervised learning}

The first thing to ask about a learning problem is \emph{what the data looks
like}. This largely determines which family of methods is relevant.

\paragraph{Supervised learning.}
Here each data point comes as a \emph{pair} $(x, y)$: an \emph{input} (or
feature vector) $x$ together with a \emph{label} (or response) $y$ that we would
like to predict. The training set is a list
\[
(x_1, y_1), \dots, (x_n, y_n),
\]
and the goal is to learn a function $f$ such that $f(x) \approx y$, so that we
can predict the label of a \emph{new} input. The name ``supervised'' reflects
the fact that, for every training input, a supervisor has provided the correct
answer $y$. Predicting the price of a flat from its characteristics, or deciding
whether an email is spam, are typical examples.

\paragraph{Unsupervised learning.}
Now the data consists of \emph{inputs only},
\[
x_1, \dots, x_n,
\]
with no labels attached. Since there is nothing specific to predict, the goal
shifts to \emph{discovering structure} in the data on its own. Typical tasks
include \emph{clustering} (grouping similar points together) and \emph{dimensionality
reduction} (finding a few directions that capture most of the variability). We will see these in \Cref{sec:unsupervised}, through the $K$-means algorithm and principal component
analysis respectively.

\paragraph{Self-supervised learning.}
A particularly powerful idea, and the engine behind most modern large-scale
models, is to \emph{manufacture the labels from the unlabelled data itself}.
In \emph{self-supervised} learning one starts from inputs only and creates an
artificial supervised problem by hiding part
of each data point and asking the model to predict it from the rest. No human
annotation is required, yet the resulting task looks supervised. The typical
example is language modelling: take a sentence, hide a word (or simply ask for
the next one), and predict it from its surrounding context; repeating this over
enormous text corpora is how large language models are trained. The same recipe
applies to images (predict missing patches) or audio. Having learned rich
representations this way, the model can then be adapted to an actual downstream
task using comparatively few labels.

\begin{remark}[Other paradigms]
Supervised and unsupervised learning are the two pillars of this course, but
they do not exhaust machine learning. In \emph{reinforcement learning}, for
instance, an agent interacts with an environment and learns from \emph{rewards}
rather than from labelled examples: think of a program learning to play a
game by trial and error. We will not treat it here.
\end{remark}

\subsection{Supervised learning: regression vs.\ classification}

Within supervised learning, problems are further split according to the
\emph{nature of the label} $y$ we want to predict.

\paragraph{Regression.}
When the label is a \emph{real-valued} (continuous) quantity, $y \in \R$, we
speak of a \emph{regression} problem. Predicting a price, a temperature, or a
person's age from a photograph are regression tasks. This is the setting we
develop in detail in the next section.

\paragraph{Classification.}
When the label takes its values in a \emph{finite set of categories}, we speak
of a \emph{classification} problem. With two categories,
$y \in \{0, 1\}$ (or $\{-1, +1\}$), the problem is called \emph{binary}
classification, e.g.\ spam vs.\ not spam, healthy vs.\ sick. With more than
two categories, $y \in \{1, \dots, K\}$, it is \emph{multi-class}, e.g.\
recognising which digit, from $0$ to $9$, is written in an image.

\begin{remark}[Same recipe, different label]
Regression and classification share the very same blueprint: a class of
candidate functions, a way to measure the quality of a prediction, and an
algorithm that searches for a good function. What changes is mostly the type of
$y$ and, consequently, the loss used to compare a prediction with the truth.
Mastering regression therefore gives us most of the tools needed for
classification as well.
\end{remark}
